\documentclass[10pt,a4paper]{amsart}
\usepackage[margin=19mm]{geometry}
\usepackage{amsmath,amssymb,mathtools}
\usepackage[scr=rsfs,cal=euler]{mathalfa}
\usepackage{xfrac}
\usepackage[unicode,hidelinks,bookmarks=false]{hyperref}
\newcommand{\ie}{i.e.\ }
\newcommand{\cf}{cf.\ }
\newcommand{\resp}{respectively\ }
\newcommand{\Subsection}{\subsection}
\providecommand{\nobreakdash}{\nobreak\hbox{-}}
\setlength{\emergencystretch}{2em}
\multlinegap0pt
\usepackage{amssymb}
\usepackage{tikz}
\newtheorem{lemma}{Lemma}
\newtheorem{theorem}{Theorem}
\title{Meta-automatic Sequences}
\author{John M. Campbell, Benoit Cloitre}
\date{Source: arXiv:2602.23395v2 (CC BY 4.0)}
\begin{document}
\maketitle

\noindent\emph{Source and licence.} Meta-automatic Sequences. Version arXiv:2602.23395v2 (CC BY 4.0), distributed by arXiv under the Creative Commons Attribution 4.0 International licence, http://creativecommons.org/licenses/by/4.0/, as recorded in the arXiv licence metadata for this exact version and shown on the abstract page https://arxiv.org/abs/2602.23395v2. Full licence notice: \texttt{licenses/arxiv-2602.23395-CC-BY-4.0.txt}.\par\smallskip

\noindent\textit{Editorial scope.} Counted unit: the article's theorem thm:M2-affine (source0.tex lines 613--622). The article's complete original proof of it is delivered with it (source0.tex lines 624--675, one \texttt{proof} environment of 52 lines). No mathematics is written, repaired or re-derived: the statement and the proof are the article's own lines, in the article's own notation, and no retained line is substituted or re-flowed.  Retained uncounted as the source-local context the proof needs: lines 243--245; lines 604--606.  Nothing was omitted from any retained span, so the package carries no omission record.  No retained line contains a citation, so the wrapper carries no bibliography and no bibliography entry.  No retained line contains a figure environment, an image-inclusion command or drawing code, so no image is delivered and the package declares no asset dependency.  Editorial formatting only, in addition: an AMS \texttt{amsart} preamble that restates the article's own declarations and macros used by the retained lines, namely the environments \texttt{lemma}, \texttt{theorem} on separate counters, together with the standard packages those lines need.  The retained environments print in this document as Lemma 1, Theorem 1 in order of appearance, whereas the article prints the same items with the numbering of its own full document.  The complete native source of the article as distributed in the arXiv e-print for this exact version is bundled unchanged as \texttt{sources/195397/source0.tex}.
\begin{lemma}\label{lem:balance-xor} 
 If $a$ is balanced, then $$ a(2n+a(n)) = a(2n)\oplus a(n) \quad\text{and}\quad a(2n+1-a(n)) = a(2n+1)\oplus a(n), $$ for all $n\ge 0$. 
\end{lemma}

\begin{lemma}\label{lem:M2-recover} 
 If $s(n)=(x,y)$, then $$\mathcal{M}_{2}(2 n) = x\oplus y, \qquad \mathcal{M}_{2}(2n+1)=x\oplus y\oplus 1.$$
\end{lemma}

\begin{theorem}\label{thm:M2-affine} 
Writing $s(n)=(x,y)$, we have
\begin{align*}
 s(4n) &= (y,\ x\oplus y),\\
 s(4n+1) &= (y\oplus 1,\ x\oplus y),\\
 s(4n+2) &= (y\oplus 1,\ x\oplus y\oplus 1),\\
 s(4n+3) &= (y,\ x\oplus y\oplus 1).
\end{align*}
Consequently $\mathcal{M}_{2}$ is $4$-automatic with at most $4$ states.
\end{theorem}

\begin{proof}
 The case $n=0$ is immediate from the initial values $\mathcal{M}_{2}(0)=0$
 and $\mathcal{M}_{2}(1)=1$. We may therefore assume $n\ge 1$ whenever a rule for $4n$ or
 $4n+1$ is used.
 Write $s(n) = (x,y)$ so that $\mathcal{M}_{2}(2n) = x \oplus y$ and $\mathcal{M}_{2}(2n+1) = x \oplus y \oplus 1$.

\textbf{Case $r = 0$.}
 The first coordinate is given by $\mathcal{M}_{2}(4n) = \mathcal{M}_{2}(2n + x)$. By Lemma~\ref{lem:M2-recover}, we have that
 $\mathcal{M}_{2}(2n) = x \oplus y$ and $\mathcal{M}_{2}(2n+1) = x \oplus y \oplus 1$. In both cases, for $x = 0$ or $x = 1$, a direct 
 check gives that $\mathcal{M}_{2}(2n + x) = y$. For the second coordinate $\mathcal{M}_{2}(8n + y)$, we see that if $y = 0$, then 
 $\mathcal{M}_{2}(8n) = \mathcal{M}_{2}(4(2n)) = \mathcal{M}_{2}(2 \cdot 2n + \mathcal{M}_{2}(2n))$, and, since $\mathcal{M}_{2}(2n) = x 
 \oplus y = x$, this equals $\mathcal{M}_{2}(4n + x)$. Simplifying both cases gives~$x$. If $y = 1$, then $\mathcal{M}_{2}(8n+1) = 1 
 - \mathcal{M}_{2}(8n)$, and, by the same analysis as before, we get $x \oplus 1$. In both sub-cases the second coordinate equals 
 $x \oplus y$. Therefore, we obtain 
$$ s(4n) = (y,\; x \oplus y). $$

\textbf{Case $r = 1$.} The first coordinate is given by $\mathcal{M}_{2}(4n+1) = 1 - \mathcal{M}_{2}(4n) = 1 - y = y \oplus 1$.

 For the second coordinate of $s(4n+1)$, by definition, this is $\mathcal{M}_{2}(2(4n+1) + \mathcal{M}_{2}(4n+1)) = 
 \mathcal{M}_{2}(8n + 2 + (y \oplus 1))$. When $y = 0$, we have $\mathcal{M}_{2}(8n+3) = \mathcal{M}_{2}(4(2n)+3) = 1 - 
 \mathcal{M}_{2}(2(2n)+1-\mathcal{M}_{2}(2n))$. Since $\mathcal{M}_{2}(2n) = x$, this equals $1 - \mathcal{M}_{2}(4 n + 1 - 
 x) = x \oplus y$ for both values of~$x$. When $y = 1$, we have $\mathcal{M}_{2}(8n+2) = \mathcal{M}_{2}(2(2n) + 1 - 
 \mathcal{M}_{2}(2n))$. Since $\mathcal{M}_{2}(2n) = x \oplus 1$, this equals $\mathcal{M}_{2}(4n+x) = x \oplus y$ for both values of 
 $x$. The second coordinate is therefore $x \oplus y$ in both cases. Therefore
\[
 s(4n+1) = (y \oplus 1,\; x \oplus y).
\]

\textbf{Case $r = 2$.} For the first coordinate $\mathcal{M}_{2}(4n+2) = \mathcal{M}_{2}(2n + 1 - x)$, Lemma~\ref{lem:balance-xor} gives $\mathcal{M}_{2}(2n+1) \oplus x = (x \oplus y \oplus 1) \oplus x = y \oplus 1$.

 For the second coordinate of $s(4n+2)$, we find that $\mathcal{M}_{2}(2(4n+2)+\mathcal{M}_{2}(4n+2)) = \mathcal{M}_{2}(8 n + 4 + 
 (y\oplus 1))$. Since $\mathcal{M}_{2}(2n+1) = x \oplus 1$, we have $\mathcal{M}_{2}(8 n + 4) = \mathcal{M}_{2}(4n+3) = 1-(y \oplus 
 1)$ for both values of~$x$. When $y=0$, this gives $\mathcal{M}_{2}(8n+4)=0$, so $\mathcal{M}_{2}(8n+5)=1=x\oplus y\oplus 1$.
When $y=1$, since $\mathcal{M}_{2}(2n+1)=x$, we get $\mathcal{M}_{2}(8n+4) = \mathcal{M}_{2}(4n+2+x)=x\oplus y\oplus 1$ by a 
 direct case check. In both cases the second coordinate is $x \oplus y \oplus 1$. It then follows that 
\[
 s(4n+2) = (y \oplus 1,\; x \oplus y \oplus 1).
\]

\textbf{Case $r = 3$.}
 The first coordinate is given by $\mathcal{M}_{2}(4n+3) = 1 - \mathcal{M}_{2}(4n+2) = 1 - (y\oplus 1) = y$.

 For the second coordinate of $s(4n+3)$, we have that $\mathcal{M}_{2}(2(4n+3)+\mathcal{M}_{2}(4n+3)) = \mathcal{M}_{2}(8 n + 6 + 
 y)$. Since $\mathcal{M}_{2}(2n+1) = x \oplus 1$, the identity $\mathcal{M}_{2}(8n+6) = \mathcal{M}_{2}(2(2n+1) + 1 - 
 \mathcal{M}_{2}(2n+1))$ gives $\mathcal{M}_{2}(8n+6) = x \oplus y \oplus 1$ by a direct two-case check on~$x$; this handles $y = 
 0$. For $y=1$, since $\mathcal{M}_{2}(2n+1)=x$, the same formula yields $1-\mathcal{M}_{2}(4n+2+x) = x\oplus y\oplus 1$. The
 second coordinate is $x \oplus y \oplus 1$ in both cases. Consequently, we have that 
$$ s(4n+3) = (y,\; x \oplus y \oplus 1).
$$

Each transition has the form $(x,y) \mapsto A\binom{x}{y} + \mathbf{v}(d)$ over $\mathbb{F}_2$, confirming the affine structure.
\end{proof}

\end{document}
